\section*{附录 · 术语速查表}
\addcontentsline{toc}{section}{附录 · 术语速查表}
\markboth{附录 · 术语速查表}{附录 · 术语速查表}

正文对每个术语都有解释（少数术语首次露面时只作预告、到对应模块才详细展开）。
此表按各术语\textbf{主要讲解位置}的先后排列，供汇报前快速复习。

\begin{center}
\small
\begin{tabular}{p{3.4cm}p{11.2cm}}
\toprule
术语 & 一句话解释 \\
\midrule
\multicolumn{2}{l}{\emph{—— 模块 1 ——}} \\
LiDAR（激光雷达） & 飞机向地面发激光、按反射时间测距的技术，产出即点云 \\
点云 & LiDAR 测得的大量三维坐标点的集合，像``由点组成的三维照片'' \\
高程 & 相对统一高度基准面的高度（不是``离地面多高''） \\
ftUS & 美制测量英尺，$\approx0.3048$ 米；美国官方测绘数据的惯例单位 \\
瓦片（tile） & 大区域数据切成的分块文件，本项目共 19 块 \\
LAZ & 点云的压缩文件格式（LAS 是标准点云格式，LAZ 是其压缩版） \\
语义分割 & 给每个点贴类别标签（地面 2/植被 5/建筑 6/其他 1），回答``是什么'' \\
实例分割 / 单体化 & 给每个建筑点一个楼编号（instance\_id），回答``是哪一栋'' \\
欠分割 / 过分割 & 该分开的没分开（楼粘成一条）/ 不该分的被拆碎（一栋变几块） \\
Gate A & 实例级噪声闸门：规则捞候选 + 人工定罪 + 原因码存档（模块 5--6 详述） \\
\multicolumn{2}{l}{\emph{—— 模块 2 ——}} \\
CloudCompare & 免费点云查看软件，本项目每步人工抽查的工具 \\
footprint（建筑底图） & 现成的每栋楼平面轮廓多边形；本项目不用作输入，留作最终评价基准 \\
\multicolumn{2}{l}{\emph{—— 模块 3 ——}} \\
栅格化 & 把连续空间切成规则方格并统计每格内容（本项目格边长 0.5 m） \\
占据格 & 落入 $\geq2$ 个建筑点的格子；全部占据格构成黑白``占据图'' \\
连通域 & 黑白图上互相挨着的黑格归成的``块''；8-连通 = 斜角也算挨着 \\
形态学操作 & 对黑白图的修形操作家族（膨胀、腐蚀、闭运算等） \\
膨胀 / 腐蚀 & 黑色区域向外长一圈 / 向内削一圈（厚度由结构元定） \\
结构元 & 规定每次扩/削多远的小模板（如 $3\times3$ = 周围一圈） \\
闭运算 & 先膨胀再腐蚀；会把楼间窄缝焊死——v1 的事故根源，v2 已弃用 \\
填洞（fill holes） & 只填被完全包围的内部空洞、不碰向外开口的缝——v2 的补洞方式 \\
\multicolumn{2}{l}{\emph{—— 模块 4 ——}} \\
核（core） & 块被腐蚀 1.5 m 后剩下的 $\geq6$ m\textsuperscript{2} 主体；个数用于估计栋数，并充当分水岭水源 \\
分水岭算法 & 把图当地形灌水、两片水相遇处筑坝作为分割边界的经典算法 \\
marker & 分水岭的指定水源；本项目用``核''充当，压住过分割 \\
$z_{\max}$ 高度图 & 每格记录格内最高点高程；取负后作为分水岭的``地形'' \\
中位数 / 百分位 & 排队取正中间的值 / 排队取指定百分比位置的值（$P_5$ = 从低往高 5\% 处） \\
碎片吸收 & $<20$ m\textsuperscript{2} 的碎片并入\textbf{平面距离 3 m 内、面积 $\geq50$ m\textsuperscript{2}、中位高程差 $\leq3$ m} 的最近``锚''实例（锚 = 吸收方；只许小并大） \\
孤儿点回挂 & 因所在格不足 2 点而暂未分块的点，回挂到 1.5 m 内最近实例格 \\
防撞色 & 每个实例随机分配且保证与邻居不同的显示颜色，避免``相邻同色''误读 \\
\multicolumn{2}{l}{\emph{—— 模块 5 ——}} \\
DTM & 数字地形模型：一张``每个位置的地面高程''图，由类别 2 地面点建成 \\
HAG（离地高度） & 实例中位高程 $z_{50}$ 减周围 4 m 环带内 DTM 的 $P_5$（环带只取有实测地面点的格，不足 20 个才用填补值） \\
长宽比 AR & 实例平面形状``最长伸展 / 最宽伸展''之比（求点分布伸展最大与最小的方向再作比），细长体侦测用 \\
凸包 & ``橡皮筋箍住形状''得到的最小凸多边形 \\
充实度 solidity & 实例面积 ÷ 凸包面积；本区实测折带 0.4--0.6、建筑主体 $>0.75$（候选门槛 0.72，终裁靠人工） \\
距离变换 EDT & 算形状内每点到边界的距离；最大值 $\times2$ = 形状``最宽处宽度'' \\
伸长指标 elong & 面积 $\div$ 最宽处宽度$^2$；同面积下越``带状''越大（候选门槛 3.5），抓方正桥板的补网 \\
原因码 & 审计表中每个剔除/拆分决定的机器可读理由（如 \texttt{reject\_reviewed\_powerline}） \\
移仓存档 & 被剔除的点不删除，单独存成文件（\texttt{noise\_removed.laz} 等）可回看回滚 \\
\multicolumn{2}{l}{\emph{—— 模块 6 ——}} \\
回波（return） & 一束激光的多次反射；``中间回波''= 多次回波中\textbf{非最后一次}的回波（首次也可能是，如``第 1 次/共 2 次''），树冠的典型特征 \\
选择偏差 & 被模型错标的子集与总体特征分布不同——``植被点 94.8\% 中间回波''不适用于``被错标的那部分树点'' \\
追溯修正（retro） & 下游步骤的新证据回传、推翻 Gate A 早先判决的机制 \\
缝内地面命中 & 判连体楼的物理判据：两块屋顶间走廊内\textbf{含类别 2 地面点的 1 m 格占比} $\geq0.35$ 判两栋（激光穿过树冠缝隙打到露天地面；实心屋面下无地面回波） \\
\bottomrule
\end{tabular}
\end{center}
